\section{Related Work}
\label{sec:related_work}

\paragraph{Architectures for persistent agent memory.}
Early LLM-agent memory systems primarily preserve interaction traces and compress them into reusable textual state. MemoryBank stores long-term records and periodically forms higher-level summaries, while MemoChat organizes extended dialogue into topic-level memos~\citep{memorybank,Memochat}. MemGPT instead treats the context window as a limited resource and moves information between in-context and archival tiers~\citep{packer2023memgpt}. More recent sequence-centric approaches learn to update a bounded recurrent state, coupling memory consolidation with downstream reasoning~\citep{memagent,mem1}. A second line of work externalizes richer structure: Mem0 extracts discrete facts into vector and graph stores~\citep{chhikara2025mem0}; MemTree builds a dynamic hierarchy of increasingly abstract memories~\citep{MemTree}; Zep represents temporally evolving entities and relations in a knowledge graph~\citep{rasmussen2025zep}; and A-MEM creates linked, self-updating note networks~\citep{xu2025amem}. Composite systems such as MemoryOS, MemOS, LightMem, and SimpleMem combine multiple granularities, indices, or lifecycle policies to balance retention, retrieval quality, and cost~\citep{memoryos,MemOS,LightMem,simplemem}. These systems establish a broad design space, but each is typically evaluated with its own workloads, model backbones, prompts, and cost accounting, making reported numbers unsuitable for direct architectural comparison.

\paragraph{Taxonomies and unified analyses.}
Surveys have organized agent memory by cognitive type, temporal scope, representation, function, and update dynamics~\citep{zhang2024survey,hu2025memory}. Such taxonomies clarify what memory stores and why an agent uses it, but are primarily descriptive rather than controlled empirical comparisons. The closest study to ours is the unified framework of \citet{wu2026memoryeab}, which decomposes memory methods into extraction, management, storage, and retrieval modules, compares representative systems on LoCoMo and LongMemEval, and analyzes token cost, context scalability, evidence position, and backbone dependence. We build on this systems view but broaden both the workload and the unit of analysis. Specifically, we study five benchmark workloads spanning 11 datasets, measure evidence-level retrieval fidelity, dynamic update robustness, long-horizon stability, and wall-clock operation cost, and construct controlled variants that alter one memory module at a time. This design supports more direct attribution to representation, extraction, routing, or maintenance choices rather than only comparison among fixed end-to-end method bundles.

\paragraph{Evaluation of long-term agent memory.}
Memory-oriented benchmarks expose complementary aspects of long-horizon interaction. LoCoMo evaluates question answering, summarization, and dialogue generation over long multi-session conversations~\citep{maharana2024locomo}; LongMemEval targets information extraction, multi-session reasoning, knowledge updates, and temporal reasoning in extended user--assistant histories~\citep{wu2024longmemeval}. MemBench crosses factual and reflective memory with participation and observation settings and reports effectiveness, efficiency, and capacity~\citep{membench}. MemoryAgentBench reformulates memory evaluation as incremental interaction and covers accurate retrieval, test-time learning, long-range understanding, and conflict resolution~\citep{memoryagentbench2026}. These benchmarks are valuable workload sources, but their primary unit of analysis is the agent's aggregate capability or final answer. They do not, by themselves, reveal which internal memory operation caused a success or failure, and results across papers remain confounded by heterogeneous implementations.

\paragraph{Positioning of this work.}
We do not introduce another memory architecture or task-specific dataset. Instead, we treat agent memory as a data management pipeline comprising representation and storage, extraction, retrieval and routing, and maintenance. A shared testbed compares 12 representative systems and two reference baselines across complementary workloads, while standardized execution traces expose accuracy--cost trade-offs and module-level ablations isolate their causes. This combination connects architecture taxonomies with benchmark outcomes and enables workload-aware conclusions about when a memory design succeeds, fails, or becomes operationally impractical.
